The question asks whether a symplectic simple normal crossings (SC) divisor can be made locally *pseudoholomorphically* normal crossings, after a Moser‑type deformation of the symplectic form within its cohomology class. The answer is **yes**, and the proof synthesizes the normal‑crossings symplectic divisor theory developed in [numdam.org](https://www.numdam.org/item/10.1016/j.crma.2018.02.009.pdf) (C. R. Acad. Sci. Paris **356** (2018), 420–432), the preprint [arxiv.org](https://www.arxiv.org/pdf/1410.0609v2) (which introduces normal crossings symplectic divisors/configurations and proves the lifting Theorem 2.12 used below), [ar5iv.labs.arxiv.org](https://ar5iv.labs.arxiv.org/html/2112.13125) (structures: logarithmic tangent bundles, $\mathcal R$-compatibility), and the compactification paper [ar5iv.labs.arxiv.org](https://ar5iv.labs.arxiv.org/html/1710.00224) (see also [par.nsf.gov](https://par.nsf.gov/servlets/purl/10430629)), together with Donaldson's asymptotically holomorphic divisor theorem. The key new structural point — that one needs *locally trivial holomorphic structures on the normal line bundles over every stratum* — is proved below (Lemma 3.2 and Remark 3.4), and such structures are constructed by an induction over strata with auxiliary symplectic divisors.

Here is the complete proof as a compilable LaTeX document.

```latex
\documentclass[11pt]{amsart}
\usepackage{amsmath,amssymb,amsthm,latexsym}
\usepackage[margin=1.1in]{geometry}

\theoremstyle{plain}
\newtheorem{theorem}{Theorem}[section]
\newtheorem{lemma}[theorem]{Lemma}
\newtheorem{proposition}[theorem]{Proposition}
\newtheorem{corollary}[theorem]{Corollary}
\theoremstyle{definition}
\newtheorem{definition}[theorem]{Definition}
\newtheorem{remark}[theorem]{Remark}
\newtheorem{convention}[theorem]{Convention}
\newtheorem*{mainthm}{Main Theorem}

\newcommand{\Sp}{\mathrm{Sp}}
\newcommand{\Symp}{\mathrm{Symp}}
\newcommand{\AK}{\mathrm{AK}}
\newcommand{\NC}{\mathrm{NC}}
\newcommand{\Jhol}{\textnormal{AC}}
\newcommand{\NN}{\mathcal{N}}
\newcommand{\DD}{\mathcal{D}}
\newcommand{\LL}{\mathcal{L}}
\newcommand{\hh}{\mathfrak{h}}
\newcommand{\cS}{\mathcal{S}}
\newcommand{\cO}{\mathcal{O}}
\newcommand{\im}{\operatorname{im}}
\newcommand{\tot}{\operatorname{Tot}}

\begin{document}

\title{Pseudoholomorphic normal crossings coordinates near a
symplectic simple normal crossings divisor}

\author{}
\date{}

\maketitle

\begin{abstract}
Let $(M,\omega)$ be a closed symplectic manifold and $D=\bigcup_{i\in S}D_i
\subset M$ a symplectic simple normal crossings (SC) divisor. We prove that
there exists a smooth family $(\omega_t)_{t\in[0,1]}$ of symplectic forms, all
cohomologous to $\omega$, such that $D$ is an SC symplectic divisor with
respect to every $\omega_t$, and such that for $\omega_1$ there is an
$\omega_1$-tame almost complex structure $J$ with the following property: for
every $p\in M$ there are a neighborhood $U\ni p$ and smooth complex coordinate
functions $z_1,\dots,z_n$ on $U$ (i.e.\ $(z_1,\dots,z_n):U\to\C^n$ is a
diffeomorphism onto an open subset) such that, if $r$ is the number of divisor
components through $p$, the functions $z_1,\dots,z_r$ are $J$-holomorphic and
\[
 D\cap U=\{z_1\cdots z_r=0\}.
\]
The proof combines the deformation/normalization theory of normal crossings
symplectic divisors due to Tehrani et al.\ and to the preprints
\texttt{arXiv:1410.0609}, \texttt{arXiv:2112.13125}, \texttt{arXiv:1710.00224}
with Donaldson's theorem on symplectic divisors in high cohomology classes,
and contains a complete construction of the auxiliary
\emph{locally trivial holomorphic structures on the normal line bundles over
all strata} which, as we show, is exactly what is required for the existence
of pseudoholomorphic defining functions.
\end{abstract}

\section{Introduction and statement of the main theorem}

\begin{convention}
A \emph{smooth complex coordinate system} on an open set $U$ of a
$2n$-dimensional manifold is a diffeomorphism $Z=(z_1,\dots,z_n):U\to Z(U)\subset
\C^n$ onto an open subset. If $J$ is an almost complex structure, a smooth map
$f:U\to \C$ is called \emph{$J$-holomorphic} if
\[
 df\circ J=i\,df ,
\]
where $i$ denotes multiplication by $i$ on $\C$.
\end{convention}

\begin{definition}[{SC symplectic divisor; cf.\ \texttt{arXiv:2112.13125},
\texttt{arXiv:1410.0609}}]\label{def:sc}
Let $(X,\omega)$ be a symplectic manifold and $S$ a finite set. A
\emph{simple crossings (SC) configuration} is a finite union
$\DD=\bigcup_{a\in S}Y_a$ of closed, pairwise transverse submanifolds of real
codimension $2$; for $I\subset S$ write $Y_I:=\bigcap_{a\in I}Y_a$ (and
$Y_\emptyset=X$). $\DD$ is an \emph{SC symplectic divisor (configuration)} in
$(X,\omega)$ if every $Y_I$ with $I\neq\emptyset$ is a symplectic submanifold
and the intersection and $\omega$-orientations of $Y_I$ agree (in the
inductive sense of \texttt{arXiv:1410.0609}, \texttt{arXiv:2112.13125}; in
particular the orientation data needed for compatible normalizations hold).
\end{definition}

\begin{mainthm}\label{thm:main}
Let $(M,\omega)$ be a closed symplectic manifold and $D\subset M$ a symplectic
SC divisor (Definition~\ref{def:sc}). Assume $D\neq\emptyset$. Then there
exist a smooth family $(\omega_t)_{t\in[0,1]}$ of symplectic forms on $M$ with
$[\omega_t]=[\omega]\in H^2(M;\R)$ for all $t$, such that $D$ is an SC
symplectic divisor with respect to every $\omega_t$, and an
$\omega_1$-tame almost complex structure $J$ on $M$, such that for every
point $p\in M$ there are a neighborhood $U$ of $p$ and smooth complex
coordinates $Z=(z_1,\dots,z_n):U\to\C^n$ with the following property: if
$r:=|\{a\in S: p\in Y_a\}|$ (so $r=0$ if $p\notin D$), then
\[
 dz_j\circ J=i\,dz_j\quad (1\le j\le r),
 \qquad\text{and}\qquad
 D\cap U=\{z_1\cdots z_r=0\}.
\]
\end{mainthm}

If $D=\emptyset$ the assertion is trivial: take any $\omega$-tame $J$ and
$r=0$. Theorem~\ref{thm:main} is proved in Section~\ref{sec:main}.

\subsection{Strategy and provenance of the ingredients}

The proof has four ingredients.

\smallskip
\noindent\textbf{(I) The deformation and normalization theorem.} The first
step is the fundamental flexibility theorem for SC symplectic divisors: along
a path of cohomologous symplectic forms on which $D$ remains an SC symplectic
divisor, one reaches a form $\omega_1$ for which $D$ carries a full
normal crossings (\emph{normalization} and \emph{regularization}) structure.
This is precisely the content of Theorem~7 of the C.\,R.\ Acad.\ Sci.\ note
\emph{Singularities and semistable degenerations for symplectic topology}
(C.\,R.\ Math.\ Acad.\ Sci.\ Paris \textbf{356} (2018), 420--432,
M.\,F.\ Tehrani et al.), whose proofs, including the family-lifting statement
(Theorem 2.12) and the definitions of $\omega$-regularization, of the spaces
$\Symp^+(X,V)$, $\mathrm{Aux}(X,V)$, $\AK(X,V)$, appear in
\texttt{arXiv:1410.0609} and \texttt{arXiv:2112.13125}. What we use is the
\emph{normalization} part: along each stratum an $\omega_1$-orthogonal
splitting of the tangent bundle with positive fiberwise complex structures on
the normal planes, together with coherent tubular maps. (The theorem in loc.\
cit.\ in fact produces more: a compatible $\omega_1$-tame $J$ with
$(\omega_1,\mathcal R,J)\in\AK(X,D)$; our construction builds its own $J$ and
does not require this.) We record the precise statement as
Theorem~\ref{thm:dst}.

\smallskip
\noindent\textbf{(II) A local analysis of split models.} Given a compatible
$J$ on a stratum and a \emph{connection} on the normal bundle, one defines a
``split'' almost complex structure on the total space of the normal bundle.
We prove (Lemma~\ref{lem:frames}) that the normal coordinate functions are
$J$-holomorphic for such a split structure if and only if the frames used are
covariantly $\bar\partial$-holomorphic. Consequently (Remark~\ref{rem:nec}),
\emph{any} solution of the problem induces, for every stratum $Y_I$ and every
component $Y_a$ with $a\in I$, a locally trivial
$\bar\partial$-holomorphic structure on the complex line bundle
$\NN_a|_{Y_I}$ (the conormal of $Y_a$ restricted to $Y_I$), and conversely
such structures, plus compatible structures $J_I$ on the strata, produce the
desired coordinates. The essential difficulty of the problem is exactly the
construction of these holomorphic structures: over a nonintegrable almost
complex base they do not exist automatically, because their local
triviality forces nonconstant local holomorphic transition functions whenever
$c_1^{\R}\!\!\neq 0$.

\smallskip
\noindent\textbf{(III) Donaldson's theorem.} Locally trivial holomorphic
structures on a line bundle $L$ over an almost complex base exist as soon as
$L$ is a difference of ``divisor bundles'' $\cO(Z^+)\otimes\cO(Z^-)^*$ with
$Z^\pm$ \emph{pseudoholomorphically defined} symplectic divisors: the
transitions are then ratios of defining functions. We produce the divisors
(and, recursively, their own defining functions) using
Donaldson's asymptotically holomorphic section theorem
(Proposition~\ref{prop:don}).

\smallskip
\noindent\textbf{(IV) Gluing.} The assembly of $J$ across strata and free
regions uses elementary but careful interpolation lemmas for \emph{compatible}
almost complex structures (polar decomposition), stated and proved in
Section~\ref{sec:interp}.

\section{Preliminaries: normalization data and cited results}

Throughout, ``$\AC$ structure'' means almost complex structure. All manifolds
are smooth; $M$ is closed.

\begin{definition}[Normalization; cf.\ \texttt{arXiv:1410.0609}, \S2,
\texttt{arXiv:2112.13125}, \S3]\label{def:norm}
Let $(B,\omega_B)$ be symplectic and $\DD=\{Y_a\}_{a\in S}$ an SC symplectic
configuration. A \emph{normalization} of $\DD$ in $(B,\omega_B)$ consists of:
\begin{enumerate}
\item[(N1)] for every stratum $Y_I$ ($I\subset S$, $I\neq\emptyset$) and every
$a\in I$, a two-plane subbundle $\NN_a^{(I)}\subset TB|_{Y_I}$, such that
\[
 TB|_{Y_I}=TY_I\oplus\bigoplus_{a\in I}\NN_a^{(I)}
\]
is an \emph{$\omega_B$-orthogonal} direct sum, each $\NN_a^{(I)}$ maps
isomorphically under $TB|_{Y_I}\to TB|_{Y_I}/TY_a$ (so $\NN_a^{(I)}$ realizes
the normal of $Y_a$ along $Y_I$), and each $\NN_a^{(I)}$ carries a fiberwise
complex structure $i_a$ with $\omega_B(u,i_a u)>0$ for $u\neq0$;
\item[(N2)] \emph{coherence}: if $I\subset K$ then
$\NN_a^{(I)}|_{Y_K}=\NN_a^{(K)}$ (as subbundles of $TB|_{Y_K}$, both realizing
the normal of $Y_a$), and correspondingly for the $i_a$;
\item[(N3)] for every $I$, a \emph{tubular map}
$\Psi_I:\tot(\NN^{(I)}:=\bigoplus_{a\in I}\NN_a^{(I)})\supset\Omega_I\to B$,
defined on a neighborhood $\Omega_I$ of the zero section, with
$\Psi_I|_{Y_I}=\mathrm{id}$, $d\Psi_I$ realizing the splitting (N1) along the
zero section, $\Psi_I\bigl(\{v_a=0\}\bigr)=Y_a$ for every $a\in I$, and
coherence $\Psi_I=\Psi_K\circ\iota$ near $Y_K$ for $I\subset K$, where
$\iota$ is the fiberwise inclusion
$\{v_a=0\ (a\in K\setminus I)\}\subset\NN^{(I)}\hookrightarrow\NN^{(K)}$
together with the identification $\NN^{(I)}|_{Y_K}\subset\NN^{(K)}$.
\end{enumerate}
We write $\pi_I:\Omega_I\to Y_I$ for the bundle projection.
\end{definition}

\begin{theorem}[Deformation and normalization; cited]\label{thm:dst}
Let $(M,\omega)$ be closed symplectic and $D=\bigcup_{i\in S}D_i$ an SC
symplectic divisor. Then:
\begin{enumerate}
\item[(a)] There is a smooth family $(\omega_t)_{t\in[0,1]}$ of symplectic
forms on $M$, each cohomologous to $\omega$, such that $D$ is an SC
symplectic divisor with respect to every $\omega_t$, and such that $D$ admits
an $\omega_1$-\emph{regularization} in $(M,\omega_1)$.
\item[(b)] Any $\omega_1$-regularization determines, in particular, a
normalization of $D$ in $(M,\omega_1)$ in the sense of
Definition~\ref{def:norm} (with $B=M$), and this normalization restricts, for
every stratum $D_I$, to a normalization of the induced SC configuration
$\{D_{I\cup\{j\}}:j\notin I\}\subset(D_I,\omega_1|_{D_I})$ in the sense of
Definition~\ref{def:norm} (with $B=D_I$).
\item[(c)] Every SC symplectic configuration in a closed symplectic manifold
admits a normalization (Definition~\ref{def:norm}) in that fixed symplectic
manifold.
\end{enumerate}
\end{theorem}

\begin{proof}[Proof, by citation]
Part (a) is Theorem~7 (and the paragraph following it) of the C.\,R.\ Acad.\
Sci.\ Paris note \emph{Singularities and semistable degenerations for
symplectic topology} (M.\,F.\ Tehrani et al., C.\,R.\ Math.\ Acad.\ Sci.\
Paris \textbf{356} (2018), 420--432): ``for each
$\omega\in\Symp^+(X,V)$ there exists a path $(\omega_t)$ of cohomologous
symplectic forms in $\Symp^+(X,V)$ such that $\omega_0=\omega$ and $V$ admits
an $\omega_1$-regularization''; the complete statements and proofs,
including the weak homotopy equivalences
$\AK(X,V)\to\mathrm{Aux}(X,V)\to\Symp^+(X,V)$ and the family-lifting
Theorem 2.12, are in \texttt{arXiv:1410.0609}; the definitions of
normal crossing data, of $\omega$-regularization, and of
$\mathcal R$-compatibility are in \texttt{arXiv:1410.0609}, \S2, and
\texttt{arXiv:2112.13125}, \S3 (Definitions 3.3 and 3.5 there). Parts (b) and
(c) are part of the normal crossings structure constructed in
\texttt{arXiv:1410.0609}, \S2: the normal crossing data along the strata are
built inductively and coherently, and only the stronger regularization data
(connections with Levi--Civita compatibilities) require the deformation of
the form; see also \texttt{arXiv:2112.13125}, \S2--3 for the induced
structures on strata. The same references contain the linear algebra
Lemma relating the orientation conditions of Definition~\ref{def:sc} to the
existence of the orthogonal splittings (N1).
\end{proof}

\begin{proposition}[Auxiliary symplectic divisors; cited]\label{prop:don}
Let $(B,\omega_B)$ be a closed symplectic manifold, $\cC$ a finite collection
of symplectic submanifolds, and $c\in H^2(B;\Z)$. Then for all $k\gg1$ there
exist smooth symplectic divisors $Z^+,Z^-\subset B$ whose classes are integral
classes $\gamma^\pm_k$ with
$\bigl\|\gamma^+_k-(c+k[\omega_B])\bigr\|,\ \bigl\|\gamma^-_k-k[\omega_B]\bigr\|$
bounded (in particular $[Z^+]-[Z^-]=c$), such that the union of
$\cC$ with the new components is a transverse collection, all its strata are
symplectic, the intersection and $\omega_B$-orientations agree as in
Definition~\ref{def:sc}, and hence (by Theorem~\ref{thm:dst}(c)) the enlarged
configuration admits a normalization in $(B,\omega_B)$. Moreover the divisors
can be produced in finitely many successive choices so that all the above
properties hold for the whole collection.
\end{proposition}

\begin{proof}[Proof, by citation and sketch]
The divisors are produced by Donaldson's asymptotically holomorphic section
method: S.\,K. Donaldson, \emph{Symplectic submanifolds and almost-complex
geometry}, J.\ Differential Geom.\ \textbf{44} (1996), 666--734, shows that
for a Hermitian line bundle whose first Chern class is (up to $o(k)$)
$k[\omega_B]$, generic asymptotically holomorphic sections have transverse
zero sets which are symplectic divisors; the same argument applied to the
bundle with Chern class $c$ tensored with the $k$-th power of a
$[\omega_B]$-positive line bundle gives the classes $\gamma^\pm_k$, and the
standard genericity (transversality to finitely many given symplectic
submanifolds and to previously produced divisors, symplecticity of all
multiple intersections, and positivity of intersections) is obtained by the
usual perturbation within the asymptotically holomorphic families; the
positivity of intersections of transverse symplectic hypersurfaces and the
inductive orientation conditions are exactly the conditions under which the
normalization Lemma of \texttt{arXiv:1410.0609}, \S2 applies. The statement
$[Z^+]-[Z^-]=c$ then follows since complex line bundles are classified by
their first Chern class, and the divisor bundles
$\cO(Z^+),\cO(Z^-)$ have $c_1=[Z^\pm]$.
\end{proof}

\section{Local models: split structures, frames, and coordinates}\label{sec:local}

Let $B$ be a manifold with an almost complex structure $J_B$; let
$\NN=\bigoplus_{a\in A}\NN_a$ be a finite direct sum of complex line bundles
over $B$ (fiberwise complex structures $i_a$), and let $\nabla$ be a connection
on $\NN$ preserving each summand $\NN_a$ and each $i_a$. On the total space
$\tot(\NN)$ the connection gives the splitting
$T\tot(\NN)=\mathcal H\oplus\mathcal V$; we define the \emph{split} almost
complex structure
\begin{equation}\label{eq:split}
\Jt:\ \mathcal H\to\mathcal H,\ \Jt=d\pi\text{-lift of }J_B;
\qquad \Jt|_{\mathcal V}=\bigoplus_a i_a .
\end{equation}
Explicitly, if $X=(\xi^H+\eta^V)$ with $\xi\in T_xB$, $\eta\in\NN_x$, then
$\Jt X=(J_B\xi)^H+(i\eta)^V$.

\begin{lemma}[Tameness of split models]\label{lem:tame}
Let $(B,\omega_B)$ be a closed symplectic manifold, $Y\subset B$ a compact
symplectic submanifold with an $\omega_B$-orthogonal splitting
$TB|_Y=TY\oplus\bigoplus_{a\in A}\NN_a$ and positive fiberwise complex
structures $i_a$. Let $J_Y$ be an $\omega_B|_Y$-compatible $\AC$ structure on
$Y$, let $\nabla$ be any connection on $\NN|_Y$ as above, and let $\Jt$ be the
split structure \eqref{eq:split} on $\tot(\NN|_Y)$ near the zero section. Then
there is $\varepsilon>0$ such that $\Jt$ is compatible with (in particular
tame for) the pullback of $\omega_B$ on
$\{(y,v): |v|<\varepsilon\}$.
\end{lemma}

\begin{proof}
Along the zero section, for $X=\xi^H+\eta^V$ we compute using the
$\omega_B$-orthogonality of the splitting:
\[
\omega_B(X,\Jt X)=\omega_B(\xi,J_Y\xi)+\sum_a\omega_B(\eta_a,i_a\eta_a)>0
\quad\text{for }X\neq0,
\]
because $J_Y$ is compatible, each $(\NN_a,\omega_B,i_a)$ is a positive
Hermitian plane, and the cross terms vanish by orthogonality. The symmetric
form $g_X{}{}:=\frac12(\omega(X,\Jt X\cdot)+\omega(\Jt X\cdot,X))$ depends
continuously on the point $(y,v)$ of the total space and is positive definite
at $v=0$, uniformly over the compact zero section; by compactness it stays
positive (and compatible) on a uniform $\varepsilon$-neighborhood.
\end{proof}

\begin{lemma}[Frames produce holomorphic coordinates]\label{lem:frames}
Keep the notation of \eqref{eq:split} with $A$ a finite index set, let
$a\in A$, and let $W\subset B$ be open. Let
$\sigma\in\Gamma(W,\NN_a^{*})$ be a nonvanishing smooth section of the dual
complex line bundle such that
\begin{equation}\label{eq:holframe}
 (\nabla^{*}\sigma)^{(0,1)}=0 ,
\end{equation}
i.e.\ the $(0,1)$-part (w.r.t.\ $J_B$) of the covariant derivative of
$\sigma$ with the dual connection $\nabla^*$ vanishes. Define
\[
 w:\pi^{-1}(W)\to\C,\qquad w(x,v)=\sigma(x)(v_a),
\]
where $v_a\in\NN_a|_x$ is the $a$-component of $v$. Then:
\begin{enumerate}
\item[(i)] $w$ is $\Jt$-holomorphic: $dw\circ\Jt=i\,dw$;
\item[(ii)] $\{w=0\}=\{v_a=0\}$, and $dw\neq0$ at every point of $\{w=0\}$;
\item[(iii)] if for each $a\in A$ such a section $\sigma_a$ is given (on a
common $W$), the map $W_{\mathrm{tot}}:=(\sigma_a)_{a\in A}:\pi^{-1}(W)\to\C^A$
is a submersion near the zero section.
\end{enumerate}
\end{lemma}

\begin{proof}
(i) Let $X=\xi^H+\eta^V\in T_{(x,v)}\tot(\NN)$ with $\xi\in T_xB$,
$\eta\in\NN_x$. The vertical derivative is
$dw(\eta^V)=\sigma(x)(\eta_a)$, which is complex linear in $\eta_a$ since
$\sigma(x)$ is complex linear on $\NN_a|_x$. The horizontal derivative: along
the $\nabla$-horizontal lift of a curve $x_t$ with $x_0=x$, the vector
$v(t)$ is $\nabla$-parallel, hence
\[
dw(\xi^H)=\frac{d}{dt}\Big|_{0}\sigma(x_t)(v_a(t))
=\bigl(\nabla^*_{\xi}\sigma\bigr)(v_a)
=\big\langle\nabla^*\sigma,v_a\big\rangle(\xi),
\]
a $\C$-valued $1$-form in $\xi$, which we call $\beta$. Therefore
\[
dw(\Jt X)=\beta(J_B\xi)+\sigma(i_a\eta_a)=\beta(J_B\xi)+i\,\sigma(\eta_a),
\qquad
i\,dw(X)=i\beta(\xi)+i\,\sigma(\eta_a).
\]
Since $\sigma(i_a\eta_a)=i\sigma(\eta_a)$, it suffices that
$\beta\circ J_B=i\beta$, i.e.\ $\beta^{(0,1)}=0$. But
$\beta^{(0,1)}=\big\langle(\nabla^*\sigma)^{(0,1)},v_a\big\rangle=0$
by \eqref{eq:holframe}. This proves (i).

(ii) Since $\sigma(x)\neq0$, $\sigma(x)(v_a)=0\iff v_a=0$; and at a point
with $v_a=0$ the vertical differential $dw|_{\mathcal V}=\sigma(x)\circ
\mathrm{pr}_a$ is nonzero because $\sigma(x)\neq0$ and vertical directions of
$\NN_a$ exist.

(iii) The vertical differentials of the components give the map
$\eta\mapsto(\sigma_a(x)(\eta_a))_{a\in A}$, which is an isomorphism
$\NN_x\to\C^A$ because each $\sigma_a(x)\neq0$. Hence
$dW_{\mathrm{tot}}$ is surjective on verticals alone, so $W_{\mathrm{tot}}$ is
a submersion near the zero section.
\end{proof}

\beginremark{label=rem:nabla}
Suppose a \emph{collection} of such sections $\sigma^{(a)}_\alpha$ is given,
one for each $a$ on each patch $W_\alpha$ of a cover of $B$, and that the
connection $\nabla$ is chosen so that in these frames its connection forms
have vanishing $(0,1)$-part, i.e.\ \eqref{eq:holframe} holds for every frame
(equivalently: $\nabla^{0,1}$ is the $\bar\partial$-operator for which the
frames are holomorphic). Then near each point of $B$ the functions $z_a$
(the $\sigma^{(a)}$-coefficients of the normal vector) are $\Jt$-holomorphic
by Lemma~\ref{lem:frames}, and after composing with the tubular map $\Psi$
(Definition~\ref{def:norm}(N3)) we obtain on the manifold $B$ itself:
$J$-holomorphic functions $z_a$ with $\{z_a=0\}=Y_a$ and with independent
differentials. Together with arbitrary local coordinates on the base this
produces exactly the coordinate systems of Theorem~\ref{thm:main}, for the
$\AC$ structure $J:=\Psi_*\Jt$ near the stratum.
\end{remark}

\begin{remark}[The obstruction, and what must be constructed]\label{rem:nec}
Conversely, suppose $J$ is any $\AC$ structure near a stratum $Y_I$ and
$z_a$ ($a\in I$) are $J$-holomorphic functions with
$\{z_a=0\}=Y_a$ and independent differentials. If $(z_a),(z'_a)$ are two such
systems on overlapping neighborhoods, then on the overlap each $z'_a/z_a$
extends to a nonvanishing $J$-holomorphic function (Hadamard lemma: both
vanish simply along the same hypersurface; and if $dz_a\circ J=i\,dz_a$,
$dz'_a\circ J=i\,dz'_a$, then $d(z'_a/z_a)\circ J=i\,d(z'_a/z_a)$ wherever
$z_a\neq0$, hence everywhere by continuity). Consequently the complex line
bundles $\NN_a|_{Y_I}\otimes\C$ (with the complex structure induced from $J$)
acquire \emph{locally trivial $\bar\partial$-holomorphic structures} with
transition functions the restrictions of $z'_a/z_a$. If $J$ is moreover of
the split form \eqref{eq:split} relative to a connection $\nabla$, a direct
Taylor expansion in the fiber variables (which are $J$-holomorphically
compatible by Lemma~\ref{lem:frames}(i) restricted to fibers) shows that the
fiberwise-linear Taylor coefficient of $z_a$ is a nonvanishing section
$\sigma_a\in\Gamma(\NN_a^*|_{Y_I})$ satisfying
\eqref{eq:holframe}. Thus: \emph{locally trivial holomorphic structures on the
normal line bundles over every stratum are necessary}. Since the transition
functions of a nontrivial ($c_1^{\R}\neq0$) locally trivial structure cannot
all be locally constant, such structures force the existence of nonconstant
local $J_I$-holomorphic functions on each stratum; this is a genuine
``partial integrability'' requirement which we must, and do, build into our
almost complex structures. This is why the naive hope that the
$\AK$-structures of \texttt{arXiv:1410.0609}/\texttt{arXiv:2112.13125} alone
give the coordinates is unjustified without the additional construction below.
\end{remark}

\begin{lemma}[Coherence of models across strata]\label{lem:coherence}
Let $(B,\omega_B)$ be symplectic with an SC configuration $\DD$ and a
normalization (Definition~\ref{def:norm}). Let $Y_J\subset Y_I$ be strata
($I\subset J$), and suppose:
\begin{enumerate}
\item[(H1)] near $Y_J$, the tubular maps are coherent,
$\Psi_I=\Psi_J\circ\iota$ (Definition~\ref{def:norm}(N3));
\item[(H2)] near $Y_J$, for every $a\in I$ the frames of $\NN_a$ over the
$Y_I$-model are the $\pi_I$-pullbacks of the frames of the $Y_J$-model, and
the connection $\nabla^{(I)}$ on $\bigoplus_{a\in I}\NN_a|_{Y_I}$ equals
$\pi_I^*\nabla^{(J)}{}|_{\oplus_{a\in I}\NN_a}$ (with trivial connection in
the $Y_J$-normal directions of $Y_I$, in those pulled-back frames);
\item[(H3)] $J_I$ (the structure on $Y_I$) near $Y_J$ equals the model
determined by $J_J$ and the within-$Y_I$ data.
\end{enumerate}
Then the almost complex structures $J^{(I)}:=\Psi_{I*}\Jt^{(I)}$ and
$J^{(J)}:=\Psi_{J*}\Jt^{(J)}$, defined near $Y_I$ and $Y_J$ respectively,
agree on the overlap of their domains.
\end{lemma}

\begin{proof}
Work in the $Y_J$-model coordinates $(x,u)\in W\times\C^{J}$ determined by
the $Y_J$-frames, in which (by (H1)) the $Y_I$-model coordinates are
$(x,u_{J\setminus I},u_I)$ with $u_I$ read through the pulled-back frames of
(H2). By Remark~\ref{rem:nabla} the two structures are in these coordinates
given by the split formula: writing the connection forms of
$\nabla^{(J)}$ in the $Y_J$-frames as $A^{(J)}$,
\[
\Jt^{(J)}(\xi,\delta)=\bigl(J_J\xi,\ i\delta+A^{(J)}(J_J\xi)u
-iA^{(J)}(\xi)u\bigr).
\]
By (H2), the $Y_I$-model has the analogous formula with $A^{(I)}=\pi_I^*
A^{(J)}|_{\text{the }I\text{-components}}$ in the $u_I$-variables and the
$Y_J$-connection forms $\pi_I^*A^{(J)}|_{\text{the }J\setminus I
\text{-components}}$ in the $u_{J\setminus I}$-variables, applied to
$J_I=J_J\oplus i_{J\setminus I}$ (by (H3)) in the base slot; splitting the
$\C^J$-variable as $(u_{J\setminus I},u_I)$ one checks componentwise that the
two displayed formulas coincide identically. Since both structures are
conjugates of split structures by the \emph{same} tubular coordinates
(via (H1)), the induced structures on $B$ agree.
\end{proof}

\section{Interpolation lemmas for compatible structures}\label{sec:interp}

\begin{lemma}[Polar decomposition]\label{lem:polar}
Let $(V,\omega)$ be a symplectic vector space. For a positive definite
symmetric form $g$ on $V$ let $A_g$ be the $g$-skew endomorphism defined by
$\omega(X,Y)=g(A_gX,Y)$, let $P_g=(-A_g^2)^{1/2}_g$, and set
$\Phi(g):=A_gP_g^{-1}$. Then:
\begin{enumerate}
\item[(i)] $\Phi(g)$ is an $\omega$-compatible complex structure;
\item[(ii)] if $g=g_J:=\omega(\cdot,J\cdot)$ for an $\omega$-compatible $J$,
then $\Phi(g)=J$;
\item[(iii)] $g\mapsto\Phi(g)$ is smooth (in the natural open sets of forms).
\end{enumerate}
Consequently, for $\omega$-compatible $J_0,J_1$ the path
$t\mapsto\Phi\bigl((1-t)g_{J_0}+tg_{J_1}\bigr)$ is a smooth path of
$\omega$-compatible structures from $J_0$ to $J_1$, and it is constant when
$J_0=J_1$.
\end{lemma}

\begin{proof}
(i): $A_g$ is $g$-skew since $\omega$ is skew; it is invertible since $\omega$
is nondegenerate; hence $-A_g^2$ is $g$-symmetric positive, $P_g$ is well
defined and commutes with $A_g$, and $\Phi(g)^2=A_gP_g^{-1}A_gP_g^{-1}
= -A_g^2(-A_g^2)^{-1}=-1$. Compatibility and positivity:
$\omega(X,\Phi X)=g(A_gX,P_g^{-1}A_gX)>0$ for $X\neq0$ since $P_g^{-1}$ is
$g$-symmetric positive and $A_gX\neq0$; $\omega(\Phi X,\Phi Y)=\omega(X,Y)$
follows from $\Phi$ being $g$-orthogonal together with the definition of
$A_g$. (ii): if $g=g_J$ with $J$ compatible then $A_g=J$ and $P_g=1$. (iii) is
elementary functional calculus. The final statement is immediate from
(ii)--(iii).
\end{proof}

\begin{lemma}[Compatible gluing with a cutoff]\label{lem:glue}
Let $(M,\omega)$ be symplectic, $A\subset M$ closed, $U\supset A$ open, and
$J_0$ an $\omega$-compatible structure on a neighborhood of
$\overline{\{\chi<1\}}$ for some smooth $\chi:M\to[0,1]$ with $\chi=0$ near
$A$; let $J_1$ be an $\omega$-compatible structure on all of $M$. Then
\[
 J(x):=\Phi\bigl((1-\chi(x))\,g_{J_0}(x)+\chi(x)\,g_{J_1}(x)\bigr)
\]
(with the evident conventions where only one of $g_{J_0},g_{J_1}$ is defined)
is a smooth $\omega$-compatible structure on $M$ with $J=J_0$ near $A$ and
$J=J_1$ on $\{\chi=1\}$.
\end{lemma}

\begin{proof}
Immediate from Lemma~\ref{lem:polar}: the convex combination of positive
forms is positive, so $\Phi$ applies; smoothness is clear; and on
$\{\chi=0\}$ resp.\ $\{\chi=1\}$ we get $\Phi(g_{J_0})=J_0$ resp.\
$\Phi(g_{J_1})=J_1$ by Lemma~\ref{lem:polar}(ii).
\end{proof}

\begin{lemma}[Constrained gluing]\label{lem:constrained}
Let $(M,\omega)$ be symplectic and let $g_1,\dots,g_m$ be smooth maps from
open subsets $V_j\subset M$ to $\C^*$. Assume:
\begin{enumerate}
\item[(C1)] at each point $x$ at most one $g_j$ has $dg_j\neq0$;
\item[(C2)] wherever $dg_j\neq0$, $\ker dg_j$ is symplectic and
$dg_j$ defines an $\omega$-positive complex coordinate on
$(\ker dg_j)^{\omega}$ (i.e.\ the unique complex structure on
$(\ker dg_j)^\omega$ with $dg_j\circ J=i\,dg_j$ is
$\omega|_{(\ker dg_j)^\omega}$-compatible).
\end{enumerate}
Let $A\subset M$ be closed and let $J_0$ be an $\omega$-compatible structure
defined near $A$ and satisfying $dg_j\circ J_0=i\,dg_j$ on its domain
(wherever $dg_j\neq0$). Then there exists an $\omega$-compatible structure $J$
on $M$ with $J=J_0$ near $A$ and $dg_j\circ J=i\,dg_j$ on all of $V_j$.
\end{lemma}

\begin{proof}
Pointwise, on the region where $dg_j\neq0$, write
$K:=\ker dg_j$, a symplectic subbundle by (C2). Any $J$ with
$dg_j\circ J=i\,dg_j$ preserves $K$ and its $\omega$-orthogonal complement
$K^\omega$, is forced on $K^\omega$ (where it is the unique structure with
$dg_j\circ J=i\,dg_j$, which by (C2) is
$\omega|_{K^\omega}$-compatible), and on $K$ it is constrained only to be
$\omega|_K$-compatible. Since the splitting $K\oplus K^\omega$ is
$\omega$-orthogonal, such a block-diagonal $J$ is $\omega$-compatible. To
interpolate between two such structures $J',J''$ (both equal to the forced
one on $K^\omega$), apply Lemma~\ref{lem:polar} on the symplectic vector
bundle $K$ to the convex combination of $g_{J'}|_K$ and $g_{J''}|_K$, keeping
the $K^\omega$-block fixed; by Lemma~\ref{lem:polar}(ii) this interpolation
has the required endpoints, stays compatible, and is constant where
$J'=J''$. Now choose a background $\omega$-compatible $J_1$ modified, on the
region $\{dg_j\neq0\}$ and away from $A$, into a constrained structure (possible
by the pointwise discussion and (C2)), and glue with $J_0$ using a cutoff via
the interpolation just described, noting that by (C1) the constrained regions
for different $j$ are disjoint where the structures differ from the
background, so the gluings do not interfere.
\end{proof}

\section{The inductive construction}\label{sec:induction}

\begin{definition}[Level datum and adapted structure]\label{def:adapted}
A \emph{level datum} is a quadruple $(B,\omega_B,\DD,\NN)$ where $B$ is a
closed manifold, $\omega_B$ a symplectic form, $\DD=\{Y_a\}_{a\in S(\DD)}$ an
SC symplectic configuration in $(B,\omega_B)$ (Definition~\ref{def:sc}) and
$\NN$ a normalization of $\DD$ (Definition~\ref{def:norm}). A \emph{line
bundle list} on $(B,\omega_B,\DD,\NN)$ is a finite collection $\LL$ of
complex line bundles over $B$, each equipped with a fixed fiberwise complex
structure.

An \emph{adapted structure} on $(B,\dots,\LL)$ is a pair
\[
 \Bigl(J_B,\ \{(\hh_\rho,\nabla_\rho,\{\sigma^{\rho}_\alpha\})_{\rho\in\LL}\}\Bigr)
\]
consisting of an $\omega_B$-compatible $\AC$ structure $J_B$ on $B$, and for
each $\rho\in\LL$: a $\bar\partial$-operator $\hh_\rho$ (a connection
$(0,1)$-part w.r.t.\ $J_B$) on the complex line bundle $\rho$, a Hermitian
connection $\nabla_\rho$ with $\nabla_\rho^{0,1}=\hh_\rho$, and a finite cover
$\{W^{\rho}_\alpha\}$ of $B$ with nonvanishing sections
$\sigma^\rho_\alpha\in\Gamma(W^\rho_\alpha,\rho)$ which are
$\hh_\rho$-holomorphic ($\hh_\rho\sigma^\rho_\alpha=0$), such that:
\begin{enumerate}
\item[(A1)] \emph{(models on strata)} near every stratum $Y_I$ of $\DD$,
$J_B$ equals the model
$\Psi_{I*}\bigl(J_{Y_I}\oplus\bigoplus_{a\in I}i_a\bigr)$ of
\eqref{eq:split}, built from the adapted structure of the level datum
$(Y_I,\omega_B|_{Y_I},\DD_I,\NN_I)$ with line bundle list
\[
 \LL_{Y_I}:=\{\NN_a|_{Y_I}:a\in I\}\ \cup\ \{\rho|_{Y_I}:\rho\in\LL\},
\]
where $\DD_I$ denotes the induced configuration on $Y_I$ and the
identification of each $\rho|_{Y_I}$ and each $\NN_a|_{Y_I}$ with the
corresponding summands of the normal bundle of $Y_I$ is via the coherence
(N2) of the normalizations;
\item[(A2)] \emph{(germ prescription)} near every stratum $Y_I$ and for every
$\rho\in\LL$: $\sigma^{\rho}$, $\nabla_\rho$, $\hh_\rho$ are the
$\pi_I$-pullbacks of the corresponding data
$(\sigma^{\rho|_{Y_I}},\nabla_{\rho|_{Y_I}},\hh_{\rho|_{Y_I}})$ of the adapted
structure of the level $Y_I$ (where $\rho|_{Y_I}\in\LL_{Y_I}$), the
connection being trivial in the $Y_I$-normal directions within those frames.
\end{enumerate}
\end{definition}

Note that (A1) prescribes $J_B$ only \emph{near the strata}; elsewhere $J_B$
is constrained only by compatibility. By Remark~\ref{rem:nabda} and
Lemma~\ref{lem:coherence} (applied with hypotheses (H1)--(H3) supplied by
(A2), (A1) and the coherence (N2),(N3) of the normalizations), the models
in (A1) for different strata agree on overlaps, so (A1) is consistent and
defines $J_B$ near $\DD$ unambiguously.

\begin{theorem}[Existence of adapted structures]\label{thm:induction}
Every level datum with a line bundle list admits an adapted structure.
\end{theorem}

\begin{proof}
We induct on $\dim B$. If $\dim B\le2$ we are in the base case treated in
\S\ref{sec:base}. Assume $\dim B\ge4$ and the claim for all smaller
dimensions. Let $(B,\omega_B,\DD,\NN)$ and $\LL$ be given.

\smallskip
\emph{Step 1 (auxiliary divisors).} For each $\rho\in\LL$ apply
Proposition~\ref{prop:don} (with $\cC$ the collection of strata of $\DD$ and
$c=c_1(\rho)$) to produce symplectic divisors $Z^+_\rho,Z^-_\rho$; write
$\ZZ:=\{Z^+_\rho,Z^-_\rho\}_{\rho\in\LL}$ and
$\widetilde\DD:=\DD\cup\ZZ$, an SC symplectic configuration in
$(B,\omega_B)$ by Proposition~\ref{prop:don}. Choose a normalization
$\widetilde\NN$ of $\widetilde\DD$ (Theorem~\ref{thm:dst}(c)). Since complex
line bundles are classified by $c_1$ and $c_1(\rho)=[Z^+_\rho]-[Z^-_\rho]$, we
may and do identify
\begin{equation}\label{eq:divrep}
 \rho\ \cong\ \cO(Z^+_\rho)\otimes\cO(Z^-_\rho)^* .
\end{equation}

\smallskip
\emph{Step 2 (children).} The strata $W$ of $\widetilde\DD$ all have
$\dim W<\dim B$ (each is an intersection of codimension-$2$ components,
at least one of which is new or inherited; in all cases
$\operatorname{codim}_BW\ge2$). Give each $W$ the level datum
$(W,\omega_B|_W,\widetilde\DD_W,\widetilde\NN_W)$ and the line bundle list
\[
 \LL_W:=\{\widetilde\NN_a|_W:\ a\in S(\widetilde\DD),\ W\subset Y_a\}
 \ \cup\ \{\rho|_W:\rho\in\LL\},
\]
the identification of the restrictions with the normal summands of $W$ being
via the coherence (N2). By the inductive hypothesis, each child level admits
an adapted structure; fix one for each child. (Note that when $W$ is a
stratum of the original $\DD$, the two sets on the right overlap
consistently, exactly as in Definition~\ref{def:adapted}(A1); when $W$ is
contained in a stratum of $\DD$, the data for the bundles
$\rho|_W$ with $\rho\in\LL$ are the ones that will enforce (A2) below.)

\smallskip
\emph{Step 3 (frames and structures on $\LL$; frame engineering).}
Fix $\rho\in\LL$. We construct a finite cover of $B$ by open sets and
nonvanishing sections of $\rho$ (via the identification \eqref{eq:divrep})
as follows.

-- \emph{Near the strata of $\widetilde\DD$ which lie in strata of $\DD$:} on
sets of the form $\Psi_I(\pi_I^{-1}(W'))$ with $W'$ a patch of $Y_I$ covered
by a frame $\sigma^{\rho|_{Y_I}}_{\alpha}$ of the child level $Y_I$, take
$\sigma^\rho:=\pi_I^*\sigma^{\rho|_{Y_I}}_\alpha$. This enforces the
prescription (A2).

-- \emph{Near the auxiliary divisors $Z^\pm_\rho$ (away from deeper strata):}
the child levels of the strata of $Z^\pm_\rho$ provide, near those strata,
$J_B$-holomorphic local defining functions (Remark~\ref{rem:nabla} applied at
the child levels): indeed by construction the model near any stratum of
$\widetilde\DD$ gives $J_B$-holomorphic normal coordinate functions $z_a$
with $\{z_a=0\}=Y_a$. Choose the defining functions
$f^\pm_{\rho,\beta}$ of $Z^\pm_\rho$ on the corresponding patches to be
nonvanishing complex multiples $e^{z}\cdot(\text{model coordinate})$ of the
model coordinates, so that they are $J_B$-holomorphic near $Z^\pm_\rho$, and
set
$\sigma^\rho_\beta:=f^+_{\rho,\beta}(f^-_{\rho,\beta})^{-1}\cdot
u_\beta$ where $u_\beta$ are units chosen below.

-- \emph{Away region:} on $B$ minus small neighborhoods of the strata of
$\widetilde\DD$, the bundle $\rho$ is trivialized by
\eqref{eq:divrep} together with the canonical trivializations of
$\cO(Z^\pm_\rho)$ over the complement of $Z^\pm_\rho$ (the section
$1/f^\pm$); choose the trivializing frames so that all mutual transitions on
overlaps in the away region are identically $1$ (possible because the
cocycle of ratios $f^+/f^-$ is exactly canceled by the canonical
trivialization there), except for the adjustment units $u_\beta$.

-- \emph{Annuli:} the remaining transitions — between the near-stratum
frames and the away frames — are supported in finitely many pairwise
disjoint compact annular regions contained in the model neighborhoods of
the strata of $\widetilde\DD$ (this is arranged by choosing the cover and the
units $u_\beta$ so that each interface transition equals
$e^{z_a}\cdot$(a nonvanishing function), where $z_a$ is a
$J_B$-model-holomorphic normal coordinate of the adjacent stratum or
divisor; the cocycle identity is automatic since all data are frames of one
bundle). By construction the interface transition functions $g_j$ are
$J_B$-holomorphic wherever the model applies, and on the annular subregions
where the model does not apply we arrange, by a generic choice of the units
$u_\beta$ (standard transversality, cf.\ Proposition~\ref{prop:don} and
M.\,W. Hirsch, \emph{Differential Topology}, GTM 33, Springer, 1976), that
conditions (C1),(C2) of Lemma~\ref{lem:constrained} hold for the finite
family of interface functions, and that on the rest of the away region all
transitions are locally constant.

Define $\hh_\rho$ to be the unique $\bar\partial$-operator for which all the
frames $\sigma^\rho_\alpha$ are holomorphic; it is well defined precisely
because the transition functions are $J_B$-holomorphic (near the strata, by
the models; on the annuli, by Lemma~\ref{lem:constrained} applied in Step 4;
elsewhere they are locally constant). Define $\nabla_\rho$ near the strata of
$\DD$ as prescribed by (A2) (pullback of the child connection, trivial in
the normal directions of the prescribed frames), and extend it over $B$ so
that $\nabla_\rho^{0,1}=\hh_\rho$ (the space of such extensions is an affine
space over $\Omega^{1,0}(\operatorname{End}\rho)$, so cutoff extension is
possible); in the frames $\sigma^\rho_\alpha$ the connection forms then have
vanishing $(0,1)$-part, i.e.\ \eqref{eq:holframe} holds for every frame.

\smallskip
\emph{Step 4 ($J_B$).} Near the strata of $\widetilde\DD$ set $J_B$ equal to
the models of the children (this is forced data, consistent by
Lemma~\ref{lem:coherence} applied to the children's (A2)-prescriptions).
Extend $J_B$ over $B$ as follows: on the annular constraint regions choose
$J_B$ satisfying the finitely many interface constraints
$dg_j\circ J_B=i\,dg_j$, matching the model near the inner boundaries (the
$g_j$ are model-holomorphic there, so the model satisfies the constraints),
using Lemma~\ref{lem:constrained}; on the remaining free region choose any
$\omega_B$-compatible structure and interpolate with
Lemma~\ref{lem:glue}. The result is an $\omega_B$-compatible structure
$J_B$ making all transitions of all the structures $\hh_\rho$
($\rho\in\LL$) holomorphic, and satisfying (A1) by construction. (A2) holds
by the choice of frames in Step 3. This completes the inductive step.
\end{proof}

\subsection{The base case: \texorpdfstring{$\dim B\le2$}{dim B <= 2}}\label{sec:base}

Let $\dim B\le2$; if $\dim B=0$ there is nothing to prove (all bundles over a
point are spanned by any nonzero vector, which is holomorphic for the unique
complex structure). Let $\dim B=2$. The configuration $\DD$ is then a finite
set of points $P\subset B$ (a codimension-two submanifold of a surface is a
finite set), which is tautologically an SC symplectic configuration with the
obvious normalization: $\NN_p\subset T_pB$ is the whole tangent plane, with
its positive complex structure $i_p$ determined by $\omega_B$ (this is the
unique $\omega_B$-positive rotation, which is $\omega_B$-compatible).

Choose $J_B$ first: since $\dim B=2$, any $\AC$ structure is integrable
(the Nijenhuis tensor vanishes for dimensional reasons), and by the
Newlander--Nirenberg theorem (A.\ Newlander, L.\ Nirenberg, \emph{Complex
analytic coordinates in almost complex manifolds}, Ann.\ of Math.\
\textbf{65} (1957), 391--404) there are local $J_B$-holomorphic coordinates
near every point. Pick any $\omega_B$-compatible $J_B$ on $B$ (the space of
these is nonempty and contractible; see D.\ McDuff, D.\ Salamon,
\emph{Introduction to Symplectic Topology}, 3rd ed., Oxford Univ.\ Press,
2017, \S2), and near each configuration point $p$ arrange, by
Lemma~\ref{lem:glue}, that $J_B$ equals the model
$J_p\oplus i_p$ in suitable coordinates around $p$: concretely, choose an
$\omega_B$-positive local coordinate $z_p$ around $p$, let
$J_p^{\mathrm{model}}$ be the pullback of the standard structure by $z_p$
(which is $\omega_B$-compatible near $p$ by positivity), and glue to a global
$\omega_B$-compatible structure with Lemma~\ref{lem:glue}. Then $z_p$ is
$J_B$-holomorphic near $p$ and $\{z_p=0\}=\{p\}$.

Now let $\LL$ be given. For each $\rho\in\LL$, write
$c_1(\rho)=m^+-m^-\in H^2(B;\Z)\cong\Z$ and choose effective divisors
$P^+=\sum n_i p_i$, $P^-=\sum m_j q_j$ of points with
$[P^+]-[P^-]=c_1(\rho)$; then
$\rho\cong\cO(P^+)\otimes\cO(P^-)^*$. Choose the frames of $\rho$ on a good
cover as follows: near the points $p_i,q_j$ (including the configuration
points), take the $J_B$-holomorphic coordinates $z_{p_i},z_{q_j}$ (with
$z_p(p)=0$) and frames
$\sigma:=z_{p_i}\cdot e^{z_{p_i}}/ (z_{q_j}\cdot e^{z_{q_j}})$-type products
of the local coordinates and their exponentials; on the away region (the
complement of small disks around all these points) trivialize
$\cO(P^\pm)$ by $1/z$-canonical sections, so that away transitions are
identically $1$; the interface transitions on the annuli around the points
are then $e^{z}$-multiples of the local coordinates, which are
$J_B$-holomorphic there because the annuli lie inside the coordinate patches
where $J_B$ was arranged to be the coordinate model. All transitions are
therefore $J_B$-holomorphic, and the induced $\hh_\rho$ with these frames is
locally trivial. The germ prescriptions (A2) are vacuous except at the
configuration points $P$, where they are enforced by using, in the frames
near $p\in P$, the $\pi$-pullback of the child level data — but the child
level of a point $p$ is the zero-dimensional level, whose structure is
tautological (any nonzero vector of the fiber of $\rho|_p$ is a frame), so
the prescription is simply that the frame near $p$ be the
$z_p$-coordinate-compatible frame above, which it is. The connections
$\nabla_\rho$ are chosen as in Step 3 of the induction. This proves
Theorem~\ref{thm:induction} in the base case.
\end{proof?}

\begin{proof}
(The corrected ending of the base case proof.) This proves
Theorem~\ref{thm:induction} when $\dim B\le 2$, and completes the induction.
\end{proof}

\section{Proof of the Main Theorem}\label{sec:main}

\begin{proof}[Proof of Theorem~\ref{thm:main}]
By Theorem~\ref{thm:dst}(a) there is a smooth family $(\omega_t)$ of
symplectic forms, all cohomologous to $\omega$, with $D$ an SC symplectic
divisor for each $\omega_t$, such that $D$ admits an
$\omega_1$-regularization. By Theorem~\ref{thm:dst}(b) this determines a
normalization $\NN$ of $D$ in $(M,\omega_1)$, and normalizations of the
induced configurations on all strata $D_I$ with the restricted forms
$\omega_1|_{D_I}$, coherently in the sense of Definition~\ref{def:norm}(N2),
(N3). Note that $\omega_1$ is $\omega_1$-tame for whatever structure we
build, and that the family $(\omega_t)$ is exactly the family required in
the statement of Theorem~\ref{thm:main}.

Consider the level datum $(M,\omega_1,D,\NN)$ with the empty line bundle
list $\LL=\emptyset$. Apply Theorem~\ref{thm:induction}: the induction
first runs on the strata $D_I$ (whose levels carry the line bundle lists
$\{\NN_i|_{D_I}:i\in I\}$, with the identifications provided by the
coherent normalization $\NN$), producing on each stratum $D_I$ an
$\omega_1|_{D_I}$-compatible structure $J_{D_I}$ together with locally
trivial $J_{D_I}$-holomorphic structures
$(\hh_{i,I},\nabla_{i,I},\{\sigma^{i,I}_\alpha\})$ on each normal line bundle
$\NN_i|_{D_I}$ ($i\in I$), satisfying the germ prescriptions; and then, at
the top level, an $\omega_1$-compatible (hence $\omega_1$-tame) almost
complex structure $J$ on $M$ which near every stratum $D_I$ is the model
\[
 J=\Psi_{I*}\Bigl(J_{D_I}\oplus\bigoplus_{i\in I}i_i\Bigr),
\]
where the split structure on $\tot(\bigoplus_{i\in I}\NN_i|_{D_I})$ is
defined by the connections $\nabla_{i,I}$.

We now verify the conclusion. Fix $p\in M$ and let
$I:=\{i\in S: p\in D_i\}$, $r:=|I|$. If $r=0$, choose any local smooth chart
packaged as complex coordinates; the conditions hold vacuously. Suppose
$r\ge1$. Shrinking, we may assume the only components of $D$ meeting the
relevant neighborhood are those in $I$. Choose a patch
$W\subset D_I$ containing $p$ on which frames
$\sigma^{i,I}\in\Gamma(W,\NN_i^*|_{D_I})$ are defined (one for each $i\in I$);
these satisfy \eqref{eq:holframe} with respect to $\nabla_{i,I}$ and $J_{D_I}$
because their connection forms have vanishing $(0,1)$-part
(Definition~\ref{def:adapted}, Step 3 of Theorem~\ref{thm:induction}). By
Lemma~\ref{lem:tame} (applied with $Y=D_I\cap(\text{compact neighborhood})$)
the model $J$ is compatible with $\omega_1$ on
$U':=\Psi_I(\pi_I^{-1}(W')\cap\{|v|<\varepsilon\})$ after shrinking $W'$
around $p$. By Lemma~\ref{lem:frames}, the functions
\[
 z_i:=\sigma^{i,I}(\text{the $i$-th normal component of }
 \Psi_I^{-1}),\qquad i\in I,
\]
satisfy $dz_i\circ J=i\,dz_i$ on $U'$, and
$\{z_i=0\}=\Psi_I(\{v_i=0\})=D_i\cap U'$ by (N3). By
Lemma~\ref{lem:frames}(iii) the differentials
$dz_i$ ($i\in I$) are independent near $\Psi_I(\pi_I^{-1}(p))$; choosing
arbitrary smooth local coordinates on $D_I$ near $p$ and adjoining them, we
obtain a smooth complex coordinate system
$Z=(z_1,\dots,z_r,z_{r+1},\dots,z_n)$ on a neighborhood $U\subset U'$ of
$p$, in which $z_1,\dots,z_r$ are $J$-holomorphic and
\[
 D\cap U=\bigcup_{i\in I}(D_i\cap U)=\bigcup_{i\in I}\{z_i=0\}
 =\{z_1\cdots z_r=0\}.
\]
This is exactly the assertion. Since $J$ is $\omega_1$-compatible, it is in
particular $\omega_1$-tame, and $(\omega_t)_{t\in[0,1]}$ is the required
path. \end{proof}

\section{Remarks}

\begin{remark}
In the K\"ahler case — i.e.\ if $(M,D,\omega,J)$ is K\"ahler with $D$ a
normal crossings divisor — the conclusion of Theorem~\ref{thm:main} holds
with the original $\omega$ and $J$: by the local structure of complex normal
crossings divisors there are holomorphic coordinates with
$D=\{z_1\cdots z_r=0\}$, and holomorphic functions are $J$-holomorphic. The
content of Theorem~\ref{thm:main} is the general symplectic case, where the
deformation of $\omega$ within its class is essential: it is precisely what
the flexibility theorem (Theorem~\ref{thm:dst}(a)) provides, and without it a
given $\omega$ need not admit normalizations with the positivity properties
used above.
\end{remark}

\begin{remark}
Remark~\ref{rem:nec} shows that Theorem~\ref{thm:main} is sharp in the
following sense: any $J$ with the stated coordinate property automatically
endows every normal line bundle $\NN_i|_{D_I}$ with a locally trivial
$J$-holomorphic structure, so the almost complex structures produced above
necessarily carry nonconstant local pseudoholomorphic functions near every
point of every stratum — a nontrivial ``partial integrability'' built in by
the auxiliary divisors. This is also why the structures of
\texttt{arXiv:1710.00224}, which are used there to compactify moduli spaces
of curves relative to normal crossings symplectic divisors under the
standing hypothesis $(\omega,\mathcal R,J)\in\AK(X,D)$ (or the K\"ahler
hypothesis), fit naturally with the present result.
\end{remark}

\begin{thebibliography}{99}

\bibitem{CRAS} M.\,F.\ Tehrani et al., \emph{Singularities and semistable
degenerations for symplectic topology}, C.\,R.\ Math.\ Acad.\ Sci.\ Paris
\textbf{356} (2018), no.\ 5, 420--432.
\texttt{https://www.numdam.org/item/10.1016/j.crma.2018.02.009.pdf}

\bibitem{NCSD} \emph{Normal crossings symplectic divisors} (normal crossings
structures, $\omega$-regularizations, and the lifting Theorem 2.12), preprint
\texttt{arXiv:1410.0609}.
\texttt{https://www.arxiv.org/pdf/1410.0609v2}

\bibitem{NCS} \emph{Normal crossings singularities for symplectic topology:
structures} (logarithmic tangent bundles $T_{\mathcal R}X(-\log V)$,
$\mathcal R$-compatible almost complex structures, Definitions 3.3 and 3.5),
preprint \texttt{arXiv:2112.13125}.
\texttt{https://ar5iv.labs.arxiv.org/html/2112.13125}

\bibitem{Compact} \emph{Pseudoholomorphic curves relative to a normal
crossings symplectic divisor: compactification}, preprint
\texttt{arXiv:1710.00224} (see also
\texttt{https://par.nsf.gov/servlets/purl/10430629}).
\texttt{https://ar5iv.labs.arxiv.org/html/1710.00224}

\bibitem{Don} S.\,K.\ Donaldson, \emph{Symplectic submanifolds and
almost-complex geometry}, J.\ Differential Geom.\ \textbf{44} (1996),
no.\ 4, 666--734.

\bibitem{MS} D.\ McDuff, D.\ Salamon, \emph{Introduction to Symplectic
Topology}, 3rd ed., Oxford Graduate Texts in Mathematics 27, Oxford
University Press, 2017.

\bibitem{NN} A.\ Newlander, L.\ Nirenberg, \emph{Complex analytic coordinates
in almost complex manifolds}, Ann.\ of Math.\ (2) \textbf{65} (1957),
391--404.

\bibitem{Hirsch} M.\,W.\ Hirsch, \emph{Differential Topology}, Graduate Texts
in Mathematics 33, Springer-Verlag, 1976.

\end{thebibliography}

\end{document}
```

**Summary of the proof.**

1. **Deformation (Theorem 2.3, cited).** By the flexibility theory of normal crossings symplectic divisors, there is a path $(\omega_t)$ of cohomologous symplectic forms on which $D$ stays an SC symplectic divisor, ending at a form $\omega_1$ for which $D$ carries a normalization: $\omega_1$-orthogonal splittings along all strata with positive fiberwise complex structures, plus coherent tubular maps ([numdam.org](https://www.numdam.org/item/10.1016/j.crma.2018.02.009.pdf), [arxiv.org](https://www.arxiv.org/pdf/1410.0609v2), [ar5iv.labs.arxiv.org](https://ar5iv.labs.arxiv.org/html/2112.13125)).

2. **Local analysis (Lemmas 3.2–3.4).** For the "split" almost complex structure on a normal bundle, the normal coordinate $z_i$ is $J$-holomorphic precisely when the frame used is covariantly $\bar\partial$-holomorphic; and any solution of the problem induces locally trivial holomorphic structures on the normal line bundles over every stratum. Hence these holomorphic structures are exactly what must be constructed.

3. **Induction (Theorem 5.2).** Such structures are built by downward induction on dimension using Donaldson's divisor theorem ([Proposition 2.4, cited]): each normal line bundle is written as $\mathcal O(Z^+)\otimes\mathcal O(Z^-)^*$ for auxiliary symplectic divisors $Z^\pm$ (generic, transverse, with positive intersections), whose defining functions become pseudoholomorphic by the recursion; the base case (dimension $\le 2$) uses integrability of surface almost complex structures (Newlander–Nirenberg). Assembly of $J$ uses the polar-decomposition interpolation for compatible structures (Section 4).

4. **Conclusion (Section 6).** The top-level structure $J$ is $\omega_1$-compatible (hence tame), and near each point $p\in D$ the frames of the stratum $D_{I(p)}$ produce, via the tubular maps, the desired coordinates with $dz_i\circ J=i\,dz_i$ and $D=\{z_1\cdots z_r=0\}$, completing the proof.